% ==========================================================================
% lampiran-B.tex — dihasilkan dari lampiran/lampiran-B-glosarium.md oleh tools/lampiran2tex.py
% Boleh disunting manual; menjalankan lampiran2tex.py lagi akan menimpa.
% ==========================================================================
\chapter{Glosarium}\label{glosarium}

Daftar istilah beserta padanan Inggris dan penjelasan ringkas. Mengikuti kebijakan istilah buku: pakai kata Indonesia yang memang lazim di komunitas TI/ilmu komputer; bila tidak ada, pertahankan istilah Inggris (dicetak \emph{miring}). Definisi di sini sengaja singkat --- untuk pembahasan penuh, lihat bab yang dirujuk.

Format tiap entri: \textbf{istilah (Indonesia / Inggris)} --- penjelasan ringkas \emph{(Bab)}.

\glosletter{A}

\textbf{Agregasi berjendela / \emph{windowed aggregation}} --- meringkas banyak event menjadi
satu nilai fitur dalam rentang waktu tertentu sebelum titik acuan (mis. total
transaksi 90 hari terakhir). \emph{(Bab 6, 10)}

\textbf{Anti-kebocoran / \emph{leakage prevention}} --- praktik memastikan tidak ada
informasi dari masa depan atau dari data uji yang bocor ke fitur saat pelatihan.
\emph{(Bab 2)}

\textbf{\emph{As-of join}} --- penggabungan tabel yang hanya mengambil informasi yang tersedia
\emph{sebelum} waktu acuan, menjaga kausalitas. \emph{(Bab 6)}

\textbf{Atribut / \emph{attribute}} --- kolom mentah pada data sebelum diolah menjadi fitur.
\emph{(Bab 1)}

\textbf{\emph{Autoencoder}} --- jaringan yang belajar memampatkan data ke representasi laten
lalu merekonstruksinya; dipakai untuk reduksi dimensi non-linear. \emph{(Bab 8)}

\textbf{Autokorelasi spasial / \emph{spatial autocorrelation}} --- kecenderungan nilai di
lokasi berdekatan lebih mirip daripada lokasi berjauhan (hukum pertama geografi
Tobler). \emph{(Bab 13)}

\glosletter{B}

\textbf{\emph{Bag-of-words} (BoW)} --- representasi teks sebagai hitungan kemunculan kata,
mengabaikan urutan. \emph{(Bab 11)}

\textbf{\emph{Baseline}} --- model atau representasi pembanding sederhana yang menjadi tolok
ukur; kenaikan skor diukur relatif terhadapnya. \emph{(Bab 9)}

\textbf{\emph{Batch}} --- sekumpulan sampel yang diproses bersama dalam satu langkah
pelatihan. \emph{(Bab 15)}

\textbf{\emph{Binning}} --- mengubah fitur numerik kontinu menjadi kategori bertingkat
(\emph{bin}); dapat menghilangkan sebagian informasi. \emph{(Bab 3)}

\glosletter{C}

\textbf{\emph{Clipping}} --- memotong nilai ekstrem ke batas atas/bawah yang ditetapkan untuk
meredam pengaruh \emph{outlier}. \emph{(Bab 3, 5)}

\textbf{\emph{Cross-fitting}} --- menghitung \emph{encoding} atau statistik pada lipatan data yang
berbeda dari yang diterapkan, mencegah kebocoran pada \emph{target encoding}. \emph{(Bab 4)}

\textbf{\emph{Cutoff} (waktu prediksi)} --- titik waktu acuan yang memisahkan informasi yang
boleh dipakai (masa lalu) dari yang harus diprediksi (masa depan). \emph{(Bab 6, 10)}

\glosletter{D}

\textbf{\emph{Data leakage}} --- masuknya informasi yang seharusnya tak tersedia saat
\emph{inference} ke dalam fitur pelatihan, membuat skor validasi terlalu optimistis.
\emph{(Bab 2)}

\textbf{\emph{Deep learning}} --- pembelajaran dengan jaringan saraf berlapis banyak yang
mampu belajar representasi bertingkat dari data. \emph{(Bab 15)}

\textbf{Deret waktu / \emph{time series}} --- data yang terurut menurut waktu, dengan
ketergantungan antar-waktu yang harus dijaga saat validasi. \emph{(Bab 10)}

\textbf{Dimensi / \emph{dimension}} --- jumlah fitur (kolom) yang merepresentasikan tiap
sampel. \emph{(Bab 8)}

\textbf{Distribusi / \emph{distribution}} --- pola sebaran nilai sebuah fitur; bentuknya
(\emph{skew}, ekor tebal) memengaruhi pilihan transformasi. \emph{(Bab 3)}

\glosletter{E}

\textbf{\emph{Embedding}} --- representasi padat berdimensi rendah yang menempatkan objek
serupa saling berdekatan dalam ruang vektor. \emph{(Bab 11, 15)}

\textbf{\emph{Encoding}} --- proses mengubah kategori menjadi angka yang dapat diolah model.
\emph{(Bab 4)}

\textbf{\emph{Entity embedding}} --- \emph{embedding} yang dipelajari untuk nilai-nilai kategori,
alternatif padat bagi \emph{one-hot} pada kardinalitas tinggi. \emph{(Bab 4, 15)}

\textbf{\emph{Epoch}} --- satu kali lintasan penuh atas seluruh data pelatihan. \emph{(Bab 15)}

\textbf{\emph{Expanding window}} --- jendela agregasi yang tumbuh dari awal riwayat sampai
titik acuan (berbeda dari \emph{rolling} yang lebarnya tetap). \emph{(Bab 10)}

\glosletter{F}

\textbf{\emph{Feature extractor}} --- model \emph{pretrained} yang lapisan-lapisannya dipakai untuk
menghasilkan fitur, tanpa dilatih ulang. \emph{(Bab 15)}

\textbf{\emph{Feature store}} --- sistem penyimpanan fitur terpusat agar konsisten antara
pelatihan dan \emph{inference}. \emph{(Bab 17)}

\textbf{\emph{Fine-tuning}} --- melanjutkan pelatihan model \emph{pretrained} pada data tugas
tertentu untuk mengadaptasinya. \emph{(Bab 15)}

\textbf{Fitur / \emph{feature}} --- besaran hasil rekayasa yang menjadi masukan model;
istilah rumah buku ini. \emph{(Bab 1)}

\textbf{Fitur turunan / \emph{derived feature}} --- fitur baru yang dibentuk dari satu atau
lebih atribut mentah (rasio, selisih, agregat). \emph{(Bab 6)}

\glosletter{G}

\textbf{\emph{Gradient}} --- arah dan besar perubahan galat terhadap parameter; dasar
optimasi model. \emph{(Bab 15)}

\textbf{\emph{Gradient boosting}} --- ansambel pohon yang dibangun bertahap; kuat pada data
tabular dan tahan terhadap skala serta \emph{outlier}. \emph{(Bab 6, 15)}

\textbf{Graf / \emph{graph}} --- data berupa simpul (\emph{node}) dan sisi (\emph{edge}) yang
merepresentasikan relasi. \emph{(Bab 13)}

\textbf{\emph{Graph neural network} (GNN)} --- jaringan yang belajar representasi simpul lewat
\emph{message passing} dari tetangga di graf. \emph{(Bab 13)}

\glosletter{H}

\textbf{\emph{Hashing (trick)}} --- memetakan kategori ke indeks tetap lewat fungsi \emph{hash},
menekan dimensi pada kardinalitas sangat tinggi dengan risiko kolisi. \emph{(Bab 4)}

\textbf{HOG (\emph{Histogram of Oriented Gradients})} --- deskriptor citra berbasis distribusi
arah gradien lokal. \emph{(Bab 12)}

\textbf{\emph{Horizon}} --- rentang waktu ke depan yang menjadi sasaran prediksi, dihitung
dari \emph{cutoff}. \emph{(Bab 10)}

\glosletter{I}

\textbf{Imputasi / \emph{imputation}} --- mengisi nilai hilang dengan estimasi (rata-rata,
median, model); harus dihitung di dalam \emph{pipeline} agar tak bocor. \emph{(Bab 5)}

\textbf{\emph{Inference}} --- tahap pemakaian model pada data baru; fitur harus tersedia pada
saat ini. \emph{(Bab 2, 17)}

\textbf{\emph{Input}} --- data yang masuk ke model; bentuknya ditentukan oleh rekayasa fitur.
\emph{(Bab 1)}

\textbf{Interpretabilitas / \emph{interpretability}} --- sejauh mana keputusan model dapat
dijelaskan; sering bergantung pada kebermaknaan fitur. \emph{(Bab 9)}

\glosletter{K}

\textbf{Kardinalitas / \emph{cardinality}} --- banyaknya nilai unik pada fitur kategorikal;
menentukan strategi \emph{encoding}. \emph{(Bab 4)}

\textbf{Kebocoran} --- lihat \textbf{\emph{Data leakage}}.

\glosletter{L}

\textbf{\emph{Lag}} --- nilai fitur atau target dari langkah waktu sebelumnya, dipakai sebagai
fitur pada deret waktu. \emph{(Bab 10)}

\textbf{\emph{Lasso} (regularisasi L1)} --- regularisasi yang menyusutkan sebagian koefisien
tepat ke nol, sekaligus menyeleksi fitur. \emph{(Bab 7)}

\textbf{\emph{Late fusion} / fusi terlambat} --- menggabungkan keluaran (skor) dari model
per-modalitas, bukan menyatukan fiturnya lebih awal. \emph{(Bab 14)}

\textbf{\emph{Machine learning}} --- pembelajaran pola dari data untuk membuat prediksi tanpa
aturan yang diprogram eksplisit. \emph{(Bab 1)}

\glosletter{M}

\textbf{MFCC (\emph{Mel-Frequency Cepstral Coefficients})} --- fitur audio ringkas yang lazim
untuk pengenalan suara. \emph{(Bab 12)}

\textbf{\emph{Missing value} / nilai hilang} --- sel data yang tidak terisi; polanya (acak
atau terstruktur) memengaruhi penanganan. \emph{(Bab 5)}

\textbf{Multimodal} --- data yang memadukan lebih dari satu modalitas (mis. teks +
citra + tabular). \emph{(Bab 14)}

\glosletter{N}

\textbf{NMF (\emph{Non-negative Matrix Factorization})} --- pemfaktoran matriks non-negatif
menjadi komponen aditif yang sering lebih mudah ditafsir. \emph{(Bab 8)}

\textbf{\emph{Noise} / \emph{derau}} --- variasi acak tak bermakna pada data; buku memakai istilah
\emph{noise}. \emph{(Bab 3)}

\glosletter{O}

\textbf{\emph{One-hot encoding}} --- merepresentasikan kategori sebagai vektor biner dengan
satu posisi bernilai 1; menjaga kategori terpisah tetapi memperlebar matriks.
\emph{(Bab 4)}

\textbf{\emph{Ordinal encoding}} --- memberi kategori kode angka berurut; hemat dimensi tetapi
memaksakan asumsi urutan. \emph{(Bab 4)}

\textbf{\emph{Outlier}} --- nilai yang menyimpang jauh dari sebaran umum; dapat mengganggu
model sensitif skala. \emph{(Bab 5)}

\textbf{\emph{Overfitting}} --- model terlalu menyesuaikan data latih hingga gagal
menggeneralisasi ke data baru. \emph{(Bab 7, 9)}

\glosletter{P}

\textbf{\emph{Pipeline}} --- rangkaian langkah praproses dan model yang dijalankan sebagai satu
kesatuan; kunci agar transformasi mengikuti pembagian data. \emph{(Bab 2)}

\textbf{Praktik pipeline yang benar / \emph{pipeline discipline}} --- kebiasaan memasang
semua transformasi di dalam \emph{pipeline} sehingga di-\emph{fit} hanya pada data latih.
\emph{(Bab 2)}

\textbf{\emph{Pretrained} (model)} --- model yang sudah dilatih pada data besar dan dipakai
kembali sebagai sumber representasi. \emph{(Bab 15)}

\textbf{\emph{Purge gap}} --- jeda yang sengaja dikosongkan antara periode latih dan uji pada
\emph{split} temporal untuk mencegah kebocoran lewat \emph{horizon}. \emph{(Bab 10)}

\glosletter{R}

\textbf{\emph{Recency}} --- fitur yang mengukur seberapa baru event terakhir relatif terhadap
\emph{cutoff}. \emph{(Bab 6)}

\textbf{Reduksi dimensi / \emph{dimensionality reduction}} --- memampatkan banyak fitur
menjadi sedikit sambil mempertahankan informasi penting. \emph{(Bab 8)}

\textbf{Regularisasi / \emph{regularization}} --- penalti pada kompleksitas model untuk
menekan \emph{overfitting}. \emph{(Bab 7)}

\textbf{Rekayasa fitur / \emph{feature engineering}} --- disiplin merancang representasi data
yang sah, bermanfaat, dan dapat digunakan kembali; tema utama buku ini. \emph{(Bab 1)}

\textbf{Representasi / \emph{representation}} --- bentuk data yang disajikan kepada model;
menentukan pola apa yang dapat dipelajari. \emph{(Bab 1)}

\textbf{Representasi yang dipelajari mesin / \emph{learned representation}} --- fitur yang
dipelajari otomatis oleh model (mis. \emph{embedding}), bukan dirancang tangan.
\emph{(Bab 15)}

\textbf{Representasi yang dirancang manusia / \emph{designed representation}} --- fitur yang
dibentuk secara sengaja oleh perancang berdasarkan pemahaman domain. \emph{(Bab 1, 6)}

\textbf{\emph{Rolling window}} --- jendela agregasi berlebar tetap yang bergeser di sepanjang
waktu. \emph{(Bab 10)}

\glosletter{S}

\textbf{Sarat informasi / \emph{information-rich}} --- sifat data/log/tabel yang memuat
banyak sinyal berguna. \emph{(Bab 1)}

\textbf{Seleksi fitur / \emph{feature selection}} --- memilih himpunan bagian fitur yang
relevan (metode \emph{filter}, \emph{wrapper}, \emph{embedded}). \emph{(Bab 7)}

\textbf{Skala / \emph{scale}} --- rentang nilai sebuah fitur; perbedaan skala antar-fitur
merusak model berbasis jarak bila tak disamakan. \emph{(Bab 3)}

\textbf{\emph{Spatial lag}} --- rata-rata berbobot nilai tetangga sebuah lokasi, dipakai
sebagai fitur spasial. \emph{(Bab 13)}

\textbf{\emph{Spectral embedding}} --- representasi simpul graf dari vektor eigen matriks graf.
\emph{(Bab 13)}

\textbf{Standardisasi / \emph{standardization}} --- penskalaan agar fitur bermean 0 dan
berdeviasi baku 1. \emph{(Bab 3)}

\glosletter{T}

\textbf{\emph{Target encoding}} --- menggantikan kategori dengan statistik target (mis.
rata-rata) yang terkait; kuat pada kardinalitas tinggi tetapi rawan bocor tanpa
\emph{cross-fitting}. \emph{(Bab 4)}

\textbf{\emph{Target} / sasaran} --- besaran yang ingin diprediksi model. \emph{(Bab 1)}

\textbf{TF-IDF (\emph{Term Frequency--Inverse Document Frequency})} --- pembobotan kata yang
menaikkan istilah khas dokumen dan menurunkan istilah umum. \emph{(Bab 11)}

\textbf{\emph{Tokenization} / tokenisasi} --- memecah teks menjadi satuan (\emph{token}) sebelum
divektorkan. \emph{(Bab 11, 15)}

\textbf{\emph{Training--inference skew} / \emph{training--serving skew}} --- ketidaksesuaian cara
fitur dihitung saat pelatihan versus saat \emph{inference}. \emph{(Bab 17)}

\textbf{Transformasi (fitur numerik)} --- pengubahan bentuk distribusi (mis. \emph{log},
\emph{Box--Cox}) agar lebih sesuai untuk model. \emph{(Bab 3)}

\glosletter{V}

\textbf{Validasi silang / \emph{cross-validation}} --- membagi data menjadi beberapa lipatan
untuk menilai generalisasi; variannya (\emph{Group}, \emph{TimeSeries}) menjaga struktur
data. \emph{(Bab 2, 9)}

\textbf{Vektor \emph{sparse} / \emph{sparse vector}} --- vektor yang sebagian besar elemennya nol,
khas representasi teks \emph{bag-of-words}/TF-IDF. \emph{(Bab 11)}

\glosletter{W}

\textbf{\emph{Winsorizing}} --- mengganti nilai ekstrem dengan persentil tertentu alih-alih
membuangnya. \emph{(Bab 5)}

\textbf{\emph{Word embedding}} --- \emph{embedding} untuk kata yang menangkap kedekatan makna
(mis. \emph{word2vec}). \emph{(Bab 11, 15)}

\begin{quote}
\textbf{Catatan penggunaan istilah.} Beberapa padanan Indonesia purist sengaja
dihindari karena tidak lazim di komunitas: \emph{derau} (dipakai: \emph{noise}), \emph{larik}
(\emph{array}), \emph{daring} (\emph{online}), \emph{masukan} untuk input model (dipakai: \emph{input}).
Kata \textbf{fitur} dipakai konsisten sebagai padanan \emph{feature}.
\end{quote}
